\chapter{Syntax}
\label{ch:syntax}

\standfirst{Smalltalk's grammar fits on a postcard. This chapter is that
postcard, expanded---and it is worth reading even if you know Smalltalk,
because three of the things here are 1983's rather than Pharo's.}

\section{The postcard}

Here is essentially all of it:

\begin{st}
| sum |                          "declare temporaries"
sum := 0.                        "assignment"
#(1 2 3) do: [:each |            "a block, passed as an argument"
    sum := sum + each].
Transcript show: sum printString; cr.   "a cascade"
^sum                             "return"
\end{st}

Six kinds of thing: literals, variables, message sends, blocks, assignment,
return. There are no keywords, no operators with special status, no
statements. \ct{ifTrue:}, \ct{do:} and \ct{whileTrue:} are messages, not
syntax, and Chapter~\ref{ch:control} shows how that can possibly work.

\section{Comments}

Double quotes, and they nest nowhere---the comment ends at the next quote. A
doubled quote inside a comment is a literal quote.

\begin{st}
"a comment"
"a comment with a "" quote in it"
\end{st}

\section{Literals}

\subsection{Numbers}

\begin{st}
42                  "SmallInteger"
-17                 "negative, the minus binds to the literal"
1000000000000       "LargePositiveInteger, automatically"
3.14                "Float"
1.5e2               "150.0"
1.5e-2              "0.015"
16rFF               "255 -- radix r digits"
2r1011              "11"
16r1F               "31"
\end{st}

Integers are arbitrary precision and promote themselves. \ct{100 factorial} is
158 digits long and correct.

\begin{stnote}[There are no scaled decimals]
Pharo's \ct{1.23s2} is not implemented, deliberately. It already parses here
as \ct{1.23 s2}---a unary message send to a Float---by exactly the same rule
that makes \ct{3 factorial} work without a space. Since \ct{Float} does not
implement \ct{s2}, you get a doesNotUnderstand and \ct{nil}, not a syntax
error. If you need exact decimals, use \ct{Fraction}.
\end{stnote}

\subsection{Characters}

A dollar sign and one character. The character after the dollar is taken
literally, including a space.

\begin{st}
$a          "the letter a"
$          "a space -- $ followed by one space"
$$          "a dollar sign"
\end{st}

Ask a character for its code with \ct{asciiValue} or \ct{asInteger}, and go
back with \ct{Character value:}.

\begin{stwarn}
\ct|$a value| does \emph{not} answer 97 here. \ct{Character} in the 1983
library implements \ct{asciiValue}, not \ct{value}; and \ct{Object>>value},
which this system adds for Pharo compatibility, answers the receiver. So
\ct|$a value| is \ct|$a|. Use \ct{asciiValue} or \ct{asInteger}.
\end{stwarn}

\subsection{Strings}

Single quotes. A doubled quote is one quote. Strings are mutable and are
arrays of characters.

\begin{st}
'hello'
'it''s'             "5 characters?  No -- 4: i t ' s"
'' size             "0"
\end{st}

\subsection{Symbols}

A hash and an identifier, a keyword sequence, or a binary selector. Symbols
are unique: two symbols spelled the same are the same object, which makes
identity comparison meaningful and hashing cheap.

\begin{st}
#foo
#at:put:            "numArgs is 2"
#+                  "numArgs is 1"
#'an odd one'       "quoted, when it is not an identifier"
\end{st}

\subsection{Literal arrays}

\ct|#(| \ldots \ct{)} builds an \ct{Array} at compile time. Everything inside
is a literal, and bare words inside become \emph{symbols}.

\begin{st}
#(1 2 3)
#(1 $a 'str' #sym (2 3))     "5 elements; the last is a nested Array"
#(foo bar)                   "two Symbols"
\end{st}

\begin{stwarn}[\ctp{\#(true false nil)} is three symbols]
In the 1983 grammar---which is this one---\ct{true}, \ct{false} and \ct{nil}
inside a literal array are ordinary bare words, and bare words become symbols.
Pharo special-cases them into the objects. Here, \ct|#(true) first class|
answers \ct{Symbol}. If you want the objects, use a dynamic array:
\ct|{true. false. nil}|.
\end{stwarn}

\subsection{Byte arrays}

\ct|#[| \ldots \ct{]}, elements 0 to 255. These are a post-Blue-Book
extension and work here, including nested inside a literal array.

\begin{st}
#[1 2 255]          "a ByteArray of 3 bytes"
\end{st}

\subsection{Dynamic arrays}

Braces build an \ct{Array} at \emph{run} time from arbitrary expressions,
separated by periods. This is the other post-Blue-Book literal form.

\begin{st}
{ 1. 2 + 3. Date today }
\end{st}

The difference from \ct|#(...)| is the whole point: a literal array holds
compile-time constants, a dynamic array holds whatever the expressions
evaluate to.

\subsection{The three pseudo-variables}

\ct{nil}, \ct{true} and \ct{false} are not literals but reserved names bound
to the sole instances of \ct{UndefinedObject}, \ct{True} and \ct{False}. You
cannot assign to them.

\section{Variables}

Names beginning with a lower-case letter are variables; names beginning with a
capital are globals, which in practice means classes and a handful of
system-wide objects like \ct{Transcript} and \ct{Smalltalk}.

Temporaries are declared between vertical bars, at the top of a method or at
the top of a block:

\begin{st}
| a b c |
\end{st}

Assignment is \ct{:=}, and it is an expression---it answers the value
assigned, so \ct{a := b := 0} works.

\begin{stnote}[Underscore assignment]
1983 source, including everything in \ct{sources/}, writes assignment as an
underscore: \ct|sum _ 0|. That was the ASCII rendering of a left-arrow glyph.
The Browser will show it to you that way. \ct{:=} is what you should type, and
both are accepted.
\end{stnote}

\section{Message sends}

There are exactly three kinds, and their precedence is fixed and simple.

\subsection{Unary}

One identifier, no arguments, binds tightest, and associates left to right.

\begin{st}
100 factorial
'hello' asUppercase reversed asSymbol
\end{st}

\subsection{Binary}

One or two non-alphabetic characters, one argument. All binary selectors have
\emph{equal} precedence and associate left to right.

\begin{st}
3 + 4 * 2           "14, not 11 -- there is no precedence among binaries"
3 + (4 * 2)         "11"
\end{st}

This is the single most common surprise in Smalltalk and it never stops being
worth stating: \textbf{arithmetic has no precedence}. Parenthesise.

Common binary selectors: \ct|+ - * / // \\ < > <= >= = ~= == ~~ , @ ->|.

\subsection{Keyword}

One or more \ct{keyword:} parts, each taking an argument. Binds loosest. A
keyword message with several parts is \emph{one} message whose selector is the
concatenation:

\begin{st}
anArray copyFrom: 2 to: 4         "one message, #copyFrom:to:"
2 raisedTo: 3 + 1                 "16 -- the binary + goes first"
5 between: 1 and: 10              "true"
\end{st}

\subsection{Precedence, in one line}

\begin{center}
unary $>$ binary $>$ keyword, and parentheses beat all three.
\end{center}

\begin{st}
x max: y factorial + 1 * 2
"reads as:   x max: (((y factorial) + 1) * 2)"
\end{st}

\section{Cascades}

A semicolon sends the next message to the \emph{same receiver} as the last
message before it. The value of a cascade is the value of the last message.

\begin{st}
Transcript show: 'a'; show: 'b'; cr
\end{st}

The receiver is whatever received \ct{show: 'a'}---that is, \ct{Transcript}.
Watch out when the first message is a keyword message with a complex
receiver:

\begin{st}
OrderedCollection new add: 1; add: 2; yourself
\end{st}

\ct{add: 1} is sent to \ct{OrderedCollection new}, so the cascade continues to
that new collection, and \ct{yourself} answers it. Without the trailing
\ct{yourself} the expression would answer \ct{2}, because that is what
\ct{add:} answers. \ct{yourself} exists for exactly this.

\section{Statements and returns}

Statements are separated by periods. A trailing period is allowed. \ct|^|
returns from the enclosing \emph{method}---see Chapter~\ref{ch:blocks} for
what that means inside a block.

A method with no explicit return answers \ct{self}.

\section{Method syntax}

A method is a pattern, optional pragmas, optional temporaries, and statements:

\begin{st}
copyFrom: start to: stop
    "Answer a copy of my elements from start to stop."
    <somePragma: 42>
    | result |
    result := self species new: stop - start + 1.
    ...
    ^result
\end{st}

The first line is the message pattern, and it is one of the three kinds:
\ct{size} (unary), \ct{+ aNumber} (binary), \ct{at: k put: v} (keyword).

By convention the first thing in the body is a comment saying what the method
answers. The Browser shows it, the debugger shows it, and the class comment
plus the method comments are the documentation.

\section{The whole grammar}

\begin{plain}
method     := pattern pragma* temporaries? pragma* statements
pattern    := unarySel | binarySel arg | (keyword arg)+
temporaries:= '|' identifier* '|'
statements := (statement '.')* (statement '.'?)? ('^' expression '.'?)?
expression := (identifier ':=')* cascade
cascade    := keywordMsg (';' message)*
keywordMsg := binaryMsg (keyword binaryMsg)*
binaryMsg  := unaryMsg (binarySel unaryMsg)*
unaryMsg   := primary identifier*
primary    := identifier | literal | block | '(' expression ')'
            | '{' (expression '.')* expression? '}'
block      := '[' (':' identifier)* ('|' temporaries?)? statements ']'
literal    := number | 'char | 'string' | #symbol | #(array)
            | #[bytes] | true | false | nil
\end{plain}

That is the language. Everything else in this book is library.
